\documentclass{beamer}

\mode<presentation>
{
  \usetheme{Singapore}      % or try Darmstadt, Madrid, Warsaw, ...
  \usecolortheme{crane} % or try albatross, beaver, crane, ...
  \usefonttheme{serif}  % or try serif, structurebold, ...
  \setbeamertemplate{navigation symbols}{}
  \setbeamertemplate{caption}[numbered]
} 

\usepackage[autostyle=true]{csquotes}
\usepackage{xeCJK}
\usepackage[T1]{fontenc}
\usepackage{xcolor}
\usepackage{setspace}
\setmainfont[AutoFakeSlant=0.2]{Charis SIL}
\setCJKmainfont[AutoFakeSlant=0.2]{Noto Serif CJK HK}
\setCJKmonofont[AutoFakeSlant=0.2]{Noto Serif CJK HK}
\usepackage{polyglossia}
\setmainlanguage{english}
\usepackage{csquotes}
\usepackage{array}
\usepackage{makecell}

\title[Pearl River Sinitic]{Yuè \& the Sinitic Languages of the Pearl River Region --- Unit 2.1}
\subtitle{13ᵗʰ EACL Summer School, Universität Stuttgart}
\author{Hilário de Sousa 蘇沙}
\institute{EHESS–CRLAO \\ {\tiny membre associé}}
\date{Tuesday 28 July 2026 \\ 丙午六月十五 \\ {\footnotesize \texttt{hilario@bambooradical.com}}}

\begin{document}

\begin{frame}
  \titlepage
\end{frame}

\section{Tones}

\begin{frame}[fragile]{Tones}
\textbf{Tone values 調值} (phonetics) vs.\\ \quad \textbf{Tone categories 調類} (phonology)\\
\bigskip
\underline{Tone values 調值}: phonetic shape of tones 
\\People commonly use Chao's tone letters or tone numerals\\
\quad {\footnotesize (Chao Yuen Ren 趙元任)}\\
e.g. Standard Mandarin\\ 
官話 /kwan˥ xwa˥˩/ \textasciitilde{} /kwan\textsuperscript{55} xwa\textsuperscript{51}/\\
你好 /ni˨˩˧꜕꜒ xao˨˩˧/ \textasciitilde{} /ni\textsuperscript{213>25} xao\textsuperscript{213}/\\
\medskip
\underline{Tone categories 調類:}\\ tones in phonological system {\footnotesize (tonemes, allotones)}\\
{\footnotesize (names like)} tones A B C D, tones 1 2 3 4, etc.\\
e.g. Mandarin 旗 /tɕʰi2/ [tɕʰi\textsuperscript{25}], Cantonese 恥 /tsʰi2/ [tsʰi\textsuperscript{25}], Taiwanese 持 /tʃʰi5/ [tʃʰi\textsuperscript{25}]\\
{\scriptsize (with numerals, be careful whether people mean tone categories or tone values)}\\
\end{frame}

\begin{frame}[fragile]{Tones}

\begin{table}
\begin{tabular}{l||c >{\centering\arraybackslash}p{0.45cm} >{\centering\arraybackslash}p{0.45cm} c||c c}
    \hline
      & \textcolor{gray!75}{*平 *A} & \multicolumn{2}{|c|}{\textcolor{gray!75}{*上 *B}} & \textcolor{gray!75}{*去 *C} & \multicolumn{2}{c}{\textcolor{gray!75}{*入 *D}}\\
     & & & & & {\scriptsize \textcolor{gray!75}{長 L}} & \multicolumn{1}{|c}{{\scriptsize \textcolor{gray!75}{短 S}}}\\ \hline
    \textcolor{gray!75}{陰 {\tiny *voiceless initial}} & 1 ˥(\~{}˥˧) & \multicolumn{2}{|c|}{2 ˨˥\~{}˧˥} & 3 ˧ & \textcolor{gray!75}{3 ˧} & \multicolumn{1}{|c}{\textcolor{gray!75}{1 ˥}}\\ \hline 
    \textcolor{gray!75}{陽 {\tiny *voiced sonorant initial}} & 4 ˨˩\~{}˩ & \multicolumn{2}{|c|}{5 ˩˧} & 6 ˨ & \multicolumn{2}{c}{\textcolor{gray!75}{6 ˨}}\\ \cline{4-4}
    \textcolor{gray!75}{陽 {\tiny *voiced obstruent initial}} & & \multicolumn{1}{|c|}{} & & & & \\ \hline
\end{tabular}
\caption{Tones in Standard Cantonese}
\end{table}

\end{frame}

\begin{frame}[fragile]{Tones}

\begin{table}
\begin{tabular}{l||c >{\centering\arraybackslash}p{0.45cm} >{\centering\arraybackslash}p{0.45cm} c||c c}
    \hline
      & \textcolor{gray!75}{*平 *A} & \multicolumn{2}{|c|}{\textcolor{gray!75}{*上 *B}} & \textcolor{gray!75}{*去 *C} & \multicolumn{2}{c}{\textcolor{gray!75}{*入 *D}}\\
     & & & & & {\scriptsize \textcolor{gray!75}{長 L}} & \multicolumn{1}{|c}{{\scriptsize \textcolor{gray!75}{短 S}}}\\ \hline
    \textcolor{gray!75}{陰 {\tiny *voiceless initial}} & 1 ˥(\~{}˥˧) & \multicolumn{2}{|c|}{2 ˨˥\~{}˧˥} & 3 ˧ & \textcolor{gray!75}{3 ˧} & \multicolumn{1}{|c}{\textcolor{gray!75}{1 ˥}}\\ \hline 
    \textcolor{gray!75}{陽 {\tiny *voiced sonorant initial}} & 4 ˨˩\~{}˩ & \multicolumn{2}{|c|}{5 ˩˧} & 6 ˨ & \multicolumn{2}{c}{\textcolor{gray!75}{6 ˨}}\\ \cline{4-4}
    \textcolor{gray!75}{陽 {\tiny *voiced obstruent initial}} & & \multicolumn{1}{|c|}{} & & & & \\ \hline
\end{tabular}
\end{table}

\begin{figure}
\flushleft
\includegraphics[width=0.55\textwidth]{cong1syu2zeoi3m4laan5do6}
{\tiny canto.hk}
\end{figure}
\medskip
倉\ 鼠\ 最\ 唔\ 懶\ 惰\ cong1 syu2 zeoi3 m4 laan5 do6\\
三\ 碗\ 半\ 牛\ 腩\ 麵\ saam1 wun2 bun3 ngau4 naam5 min6\\

\end{frame}

\begin{frame}[fragile]{Tones}

\begin{table}
\begin{tabular}{l||c >{\centering\arraybackslash}p{0.45cm} >{\centering\arraybackslash}p{0.45cm} c||c c}
    \hline
     & *A & \multicolumn{2}{|c|}{*B} & *C & \multicolumn{2}{c}{*D}\\
     & & & & & {\scriptsize L} & \multicolumn{1}{|c}{{\scriptsize S}}\\ \hline
    陰 & 1 ˥(\~{}˥˧) & \multicolumn{2}{|c|}{2 ˨˥\~{}˧˥} & 3 ˧ & 3\textsuperscript{(8)} ˧ & \multicolumn{1}{|c}{1\textsuperscript{(7)} ˥}\\ \hline 
    陽 & 4 ˨˩\~{}˩ & \multicolumn{2}{|c|}{5 ˩˧} & 6 ˨ & \multicolumn{2}{c}{6\textsuperscript{(9)} ˨}\\ \hline
\end{tabular}
\end{table}

\begin{table}
\flushleft
\begin{tabular}{l r}
\makecell[tl]{\textbf{6 tones? 9 tones?}: \\ 6 tonemes \\「九聲六調」} &
\raisebox{-\height}{\includegraphics[width=0.4\textwidth]{sixnine}}
\\
\end{tabular}
\end{table}
\bigskip
{\scriptsize \textit{Chinese linguistics traditionally counts tones in sonorant-ending syllables and obstruent-ending syllables (-p -t -k) separately}}\\

\end{frame}

\begin{frame}[fragile]{Tones}

\begin{table}
\begin{tabular}{l||c >{\centering\arraybackslash}p{0.45cm} >{\centering\arraybackslash}p{0.45cm} c||c c}
    \hline
     & *A & \multicolumn{2}{|c|}{*B} & *C & \multicolumn{2}{c}{*D}\\
     & & & & & {\scriptsize L} & \multicolumn{1}{|c}{{\scriptsize S}}\\ \hline
    陰 & 1 ˥(\~{}˥˧) & \multicolumn{2}{|c|}{2 ˨˥\~{}˧˥} & 3 ˧ & 3\textsuperscript{(8)} ˧ & \multicolumn{1}{|c}{1\textsuperscript{(7)} ˥}\\ \hline 
    陽 & 4 ˨˩\~{}˩ & \multicolumn{2}{|c|}{5 ˩˧} & 6 ˨ & \multicolumn{2}{c}{6\textsuperscript{(9)} ˨}\\ \hline
\end{tabular}
\end{table}
\bigskip
\begin{tabular}{l l}
詩 si1 \enquote*{poem} & 唧 zit1\textsuperscript{(7)} \enquote*{tickle}\\
史 si2 \enquote*{history} & \\
試 si3 \enquote*{to try} & 節 zit3\textsuperscript{(8)} \enquote*{festival}\\
時 si4 \enquote*{time} & \\
市 si5 \enquote*{city} & \\
事 si6 \enquote*{matter} & 截 zit6\textsuperscript{(9)} \enquote*{to sever}\\
\end{tabular}

\end{frame}

\begin{frame}[fragile]{Tones}

\begin{table}
\begin{tabular}{l||c >{\centering\arraybackslash}p{0.45cm} >{\centering\arraybackslash}p{0.45cm} c||c c}
    \hline
     & *A & \multicolumn{2}{|c|}{*B} & *C & \multicolumn{2}{c}{*D}\\
     & & & & & {\scriptsize L} & \multicolumn{1}{|c}{{\scriptsize S}}\\ \hline
    陰 & 1 ˥(\~{}˥˧) & \multicolumn{2}{|c|}{2 ˨˥\~{}˧˥} & 3 ˧ & 3\textsuperscript{(8)} ˧ & \multicolumn{1}{|c}{1\textsuperscript{(7)} ˥}\\ \hline 
    陽 & 4 ˨˩\~{}˩ & \multicolumn{2}{|c|}{5 ˩˧} & 6 ˨ & \multicolumn{2}{c}{6\textsuperscript{(9)} ˨}\\ \hline
\end{tabular}
\end{table}
\medskip

\begin{table}
\flushleft
\begin{tabular}{l l l l l l l l l}
\multicolumn{6}{l}{\emph{\small sonorant-ending}} & \multicolumn{3}{|l}{\emph{\small obstruent-ending}}\\ \hline
\footnotesize cong1 & \footnotesize syu2 & \footnotesize zeoi3 & \footnotesize m4 & \footnotesize laan5 & \footnotesize do6 & \footnotesize bit1\textsuperscript{(7)} & \footnotesize faat3\textsuperscript{(8)} & \footnotesize lik6\textsuperscript{(9)}\\ 
倉 & 鼠 & 最 & 唔 & 懶 & 惰 & 必 & 發 & 力\\ 
\footnotesize saam1 & \footnotesize wun2 & \footnotesize bun3 & \footnotesize ngau4 & \footnotesize naam5 & \footnotesize min6 & \footnotesize jat1\textsuperscript{(7)} & \footnotesize baak3\textsuperscript{(8)} & \footnotesize dip6\textsuperscript{(9)}\\ 
三 & 碗 & 半 & 牛 & 腩 & 麵 & 一 & 百 & 碟\\ 
%\footnotesize saam1 & \footnotesize gau2 & \footnotesize sei3 & \footnotesize ling4 & \footnotesize ng5 & \footnotesize ji6 & \footnotesize cat1\textsuperscript{(7)} & \footnotesize baat3\textsuperscript{(8)} & \footnotesize luk6\textsuperscript{(9)}\\ 
%三 & 九 & 四 & 零 & 五 & 二 & 七 & 八 & 六\\ 
\end{tabular}
\end{table}
\medskip
\vfill {\scriptsize See more examples at @語言學貓 \url{https://youtu.be/0ebW9QFlK0s?si=cqwoMFQa_BdgQG1A&t=73}\\
\emph{(From now on: tones 7 8 9 will usually be notated as tones 1 3 6, following Jyutping)}}\par
\end{frame}

\begin{frame}[fragile]{Tones}

\begin{table}
\begin{tabular}{l||c >{\centering\arraybackslash}p{0.45cm} >{\centering\arraybackslash}p{0.45cm} c||c c}
    \hline
     & *A & \multicolumn{2}{|c|}{*B} & *C & \multicolumn{2}{c}{*D}\\
     & & & & & {\scriptsize L} & \multicolumn{1}{|c}{{\scriptsize S}}\\ \hline
    陰 & 1 ˥(\~{}˥˧) & \multicolumn{2}{|c|}{2 ˨˥\~{}˧˥} & 3 ˧ & 3\textsuperscript{(8)} ˧ & \multicolumn{1}{|c}{1\textsuperscript{(7)} ˥}\\ \hline 
    陽 & 4 ˨˩\~{}˩ & \multicolumn{2}{|c|}{5 ˩˧} & 6 ˨ & \multicolumn{2}{c}{6\textsuperscript{(9)} ˨}\\ \hline
\end{tabular}
\end{table}
\bigskip
Obstruent-ending syllables, with suprafix tone 2:\\
\quad 碟 dip6*2 [tiːp˨˥] \enquote*{plate}\\
\quad 賊 caak6*2 [tsʰaːk˨˥] \enquote*{thief}\\
\medskip
Obstruent-ending syllables, tone 4 in ideophones:\\
\quad 𠽌𠽌聲 \underline{syut4} syut2 seng1 {\small sound of slurping e.g. noodle}\\
\quad 軟𢴈𢴈 jyun5 \underline{det4 det4} {\small weak and flimpsy} \textsuperscript{廣州話方言詞典}\\
\quad 濕笛笛 sap1 \underline{det4 det4} \~{} \underline{zet4 zet4} \~{} det6 det6\\ \quad {\small fully soaked and dripping} 

\end{frame}

\begin{frame}[fragile]{Tones}
Careful!\\ 
In Yuè studies, tone categories often numbered like this:\\

\begin{tabular}{c|c c c|c c}
 & A & B & C & DL & DS\\ \hline
陰 & 1 & 2 & 3 & 8 & 7\\ \hline
陽 & 4 & 5 & 6 & (10) & 9\\
\end{tabular}

\bigskip
In Mǐn studies, tone categories often numbered like this:\\
\begin{tabular}{c|c c c|c}
 & A & B & C & D\\ \hline
陰 & 1 & 2 & 3 & 4\\ \hline
陽 & 5 & 6 & 7 & 8\\
\end{tabular}

\bigskip
There is also a Pan-Sinitic system that people may use:\\
\begin{tabular}{c|c c c|c}
 & A & B & C & D\\ \hline
陰 & 1 & 3 & 5 & 7\\ \hline
陽 & 2 & 4 & 6 & 8\\
\end{tabular}


\end{frame}

\section{Suprafixes}

\subsection{Suprafixes 1}
\begin{frame}[fragile]{Suprafixes}
Suprafix: suprasegmental affix {\footnotesize (affix in the form of e.g. tone, stress)}\\
In the case of Cantonese: tonal affixes\\
\medskip
\textbf{suprafix ≠ tone-sandhi}\quad {\footnotesize although both are called 變調 in Chinese}\\
\bigskip
\underline{Tone-sandhi}:\\
\textsuperscript{(largely)} automatic \& predictable process of tones changing, \\
\quad given a particular phonological / morphosyntactic environ.\\
\medskip
e.g. Standard Mandarin tone 3 sandhi rule\\
買馬 mǎi mǎ > mái mǎ (= 埋馬 mái mǎ)\\
\smallskip
e.g. Shanghainese\\
三杯 se˥˧ pe˥˧ > se˥ pe˨˩\\
四杯 sz˧˦ pe˥˧ > sz˧ pe˦\\ 

\end{frame}

\begin{frame}[fragile]{Suprafixes}
Tone Sandhi in Cantonese?\\
\bigskip
e.g. Cheung 張 (1969)\\ 
different sandhi rules concerning ˥˧ > ˥ in Canton vs Hong Kong\\
\medskip
\emph{However, it is very hard to find people nowadays who contrast ˥˧ and ˥ consistently\\
\quad Hong Kong is nowadays basically all ˥}\\
\smallskip
Most Cantonese speakers do not have tone sandhi

\vfill \setstretch{0.65} {\tiny Cheung, Yat-Shing 張日昇. 1969. 香港粵語陰平調及變調問題 A study on the upper
even tone and changed tones in Cantonese as spoken in Hong Kong. \emph{The Journal of The Institute of Chinese Studies of The Chinese
University of Hong Kong} 香港中文大學中國文化研究所學報 2: 81–107. \url{https://www.cuhk.edu.hk/ics/journal/articles/v02p081.pdf}.} 

\end{frame}

\begin{frame}[fragile]{Suprafixes}

{\small \emph{(Suprafixes in Standard Cantonese are tone-substitutive)}}\\
Grammatical vs. lexical suprafixes:\\
\textbf{\underline{Grammatical suprafix}}:\\

\medskip
e.g. verbal:\\
\medskip
食 sik6 \enquote*{eat} > sik6*2 (= sik6 zo2 食咗  eat \textsc{pfv})\\
\smallskip
nei5 sik6*2 faan6 mei6 aa3?\\
\textsc{2sg} eat:\textsc{pfv} rice not.yet \textsc{q}\\
\enquote*{Have you eaten yet?}\\
\bigskip
氹 tam3 \enquote*{to coax}, \enquote*{to cheer s.o. up}\\
tam3*2\~{}tam3 (= tam3 jat1 tam3 氹一氹 coax:\textsc{delim})\\
\smallskip
tam3*2\~{}tam3 keoi5\\
coax:\textsc{delim} \textsc{3sg}\\
\enquote*{Coax him/her a little bit.}

\end{frame}

\begin{frame}[fragile]{Suprafixes}
\textbf{\underline{Grammatical Suprafix}}:\\
\medskip
e.g. adjectival:\\
\medskip
滑 waat6 \enquote*{smooth} > waat6*2\~{}waat6\\ 
(= waat6 jat1 waat6 滑一滑 \enquote*{extraordinarily smooth})\\
\smallskip
sik6 dou3 pei4fu1 waat6*2\~{}waat6\\
eat till skin smooth:extraordinary\\
\enquote*{Eat this and the skin becomes smooth as hell.}\\
{\tiny \texttt{https://hkppltravel.com/23966/foods/【滑溜溜豆腐花！】/}}\\

\bigskip
紅 hung4 \enquote*{red} > hung4\~{}hung4*2-dei2 \enquote*{somewhat red}\\

\end{frame}

\begin{frame}[fragile]{Suprafixes}
\textbf{\underline{Lexical (Nominal) Suprafix}}: {\footnotesize see also de Sousa (2024)}\\
\smallskip
a.k.a. 小稱變調 siu2 cing1 bin3 diu6 \enquote*{diminutive tone-change}\\
{\scriptsize However: Cantonese lexical suprafixes have been used so much that they often do not have a diminutive meaning \emph{(small / endearment meaning)}.}\\
\medskip
e.g.\\ 
\begin{table}[]
    \flushleft
    \begin{tabular}{l l}
        頭 tau4 \enquote*{head} & 糖 tong4 \enquote*{sugar} \\
        阿頭 aa3 tau4*2 \enquote*{boss, chief} & 糖 tong4*2 \enquote*{lolly, candy}
    \end{tabular}
\end{table}

\medskip
comparable with Northern Mandarin -兒 -r, e.g.\\
\begin{table}[]
    \flushleft
    \begin{tabular}{l}
        頭 tóu \enquote*{head}\\
        頭兒 tóur \enquote*{boss, chief}
    \end{tabular}
\end{table}
%{\textcolor{gray!85}{\footnotesize (somewhat comparable with\\ e.g. German diminutives:\\
%Magd \enquote*{maid, female servant}\\
%Mädchen \enquote*{girl})}}

\vfill \setstretch{0.65}{\tiny de Sousa, Hilário. 2024. Some observations on the Cantonese lexical suprafixes. \emph{Languages} 9 (10). 311. (Special issue "Typology of Chinese Languages: One Name, Many Languages", edited by Umberto Ansaldo \& Pui Yiu Szeto.) \url{https://doi.org/10.3390/languages9100311}}

\end{frame}

\begin{frame}[fragile]{Suprafixes}
Cantonese lexical suprafix:\\
Derivative morphology (word building)\\
\medskip
Usually applied to last syllable of a noun / noun stem, e.g.\\
{\footnotesize 魚 jyu4*2 \enquote*{fish}, 蝴蝶 wu4 dip6*2 \enquote*{butterfly}\\
(cf. 魚生 jyu4 saang1 \enquote*{raw fish}, 蝴蝶灣 wu4 dip6 waan1 \enquote*{butterfly bay})}\\
\medskip
Usually do not apply to other word classes, e.g.\\
{\footnotesize 鉗 kim4*2 \enquote*{pliers}, 鉗 kim4 \enquote*{to grab with pliers}}\\
{\footnotesize 碟 dip6*2 \enquote*{plate}, 碟 dip6 \textsc{clf}}\\
\medskip
Not quite predictable, e.g.\\
{\footnotesize 法蘭西 faat3 laan4*1 sai1 \enquote*{France}, 荷蘭 ho4 laan4*1 \enquote*{Holland},\\ \quad but 芬蘭 fan1 laan4 \enquote*{Finland}, 紐西蘭 nau2 sai1 laan4 \enquote*{New Zealand}\\
\smallskip
澳門 ou3 mun4*2 \enquote*{Macau},\\
\quad but 屯門 tyun4 mun4 \enquote*{Tuen Mun}, 江門 gong1 mun4 \enquote*{Kongmoon / Jiangmen}, 廈門 haa6 mun4 \enquote*{Amoy / Xiamen}}\\

\end{frame}

\begin{frame}[fragile]{Suprafixes}
\begin{table}
\begin{tabular}{l||c >{\centering\arraybackslash}p{0.45cm} >{\centering\arraybackslash}p{0.45cm} c||c c}
    \hline
     & *A & \multicolumn{2}{|c|}{*B} & *C & \multicolumn{2}{c}{*D}\\
     & & & & & {\scriptsize L} & \multicolumn{1}{|c}{{\scriptsize S}}\\ \hline
    陰 & 1 ˥(\~{}˥˧) & \multicolumn{2}{|c|}{2 ˨˥\~{}˧˥} & 3 ˧ & 3 ˧ & \multicolumn{1}{|c}{1 ˥}\\ \hline 
    陽 & 4 ˨˩\~{}˩ & \multicolumn{2}{|c|}{5 ˩˧} & 6 ˨ & \multicolumn{2}{c}{6 ˨}\\ \hline
\end{tabular}
\end{table}

Lexical suprafixes in Standard Cantonese:\\
\begin{itemize}
    \item tone 2 [˨˥] (most common)
    \item tone 1 [˥] (rarer)
    \item \textcolor{gray!75}{tone 4 [˩] (only in some informal personal names,\\ e.g. 達哥 daat6-go1*4 \enquote*{Brother Tat}; or place names, e.g. 長沙灣 coeng4 saa1 waan1*4 \enquote*{Cheung Sha Wan})}
\end{itemize}

\end{frame} 

\begin{frame}[fragile]{Suprafixes}
\begin{table}
\begin{tabular}{l||c >{\centering\arraybackslash}p{0.45cm} >{\centering\arraybackslash}p{0.45cm} c||c c}
    \hline
     & *A & \multicolumn{2}{|c|}{*B} & *C & \multicolumn{2}{c}{*D}\\
     & & & & & {\scriptsize L} & \multicolumn{1}{|c}{{\scriptsize S}}\\ \hline
    陰 & 1 ˥(\~{}˥˧) & \multicolumn{2}{|c|}{2 ˨˥\~{}˧˥} & 3 ˧ & 3 ˧ & \multicolumn{1}{|c}{1 ˥}\\ \hline 
    陽 & 4 ˨˩\~{}˩ & \multicolumn{2}{|c|}{5 ˩˧} & 6 ˨ & \multicolumn{2}{c}{6 ˨}\\ \hline
\end{tabular}
\end{table}
Suprafix tone 2 [˨˥] can be applied to\\ \quad morphemes with citation tones 3 4 5 6\\
{\scriptsize 片 pin3*2 \enquote*{film}, 房 fong4*2 \enquote*{room}, 女 neoi5*2 \enquote*{daughter}, 妹 mui6*2 \enquote*{younger sister}, 鴨 aap3*2 \enquote*{duck}, 碟 dip6*2 \enquote*{plate}}\\
\medskip
Suprafix tone 1 [˥] can be applied to\\ \quad morphemes with citation tones 2 3 4 5 6\\
{\scriptsize 自己 zi6 gei2(*1) \enquote*{self}, -派 -paai3(*1) \enquote*{faction}, 欄 laan4*1 \enquote*{animal pen}, -尾 -mei5*1 \enquote*{last}, -妹 -mui6*1 \enquote*{girl}, 掠水 loek6*1 seoi2 \enquote*{rob money}}

\end{frame}

\begin{frame}[fragile]{Suprafixes}

People usually know the citation tone, as the citation tone morpheme appears in other environments, e.g.:\\
\medskip
蛋 daan6*2 \enquote*{egg}\\
雞蛋 gai1 daan6*2 \enquote*{chicken egg}\\
雞蛋仔 [gai1 daan6*2] -zai2 ([fowl egg] \textsc{dim})\\
\medskip
雞蛋仔 gai1 daan6 zai2 \enquote*{egg waffle}\\
{\scriptsize \quad \url{https://en.wikipedia.org/wiki/Egg_waffle}}\\
蛋撻 daan6 taat1 \enquote*{egg tart}\\
{\scriptsize \quad \url{https://en.wikipedia.org/wiki/Egg_tart}}\\
蛋糕 daan6 gou1 \enquote*{cake}\\
\medskip
(蛋 daan6 is not used as an independent morpheme.)

\end{frame}

\begin{frame}[fragile]{Suprafixes}
However, the citation tone is not always known by speakers. {\scriptsize (The citation tone version of the morpheme is no longer used)}\\
\medskip
e.g. daai6 mou6 saan1 \textsuperscript{highest peak in HK} \\ 
is sometimes wrongly written as 大霧山 \enquote*{big fog mountain}\\
Should be: 大帽山 \enquote*{big hat mountain}; 帽 mou6*2 \enquote*{hat}\\
{\footnotesize (the citation tone version 帽 mou6 is no longer known)}\\
\medskip

\begin{table}
\flushleft
\begin{tabular}{l r}
\makecell[tl]橙 caang?*2 \enquote*{orange} &
\raisebox{-\height}{
\includegraphics[width=0.35\textwidth]{caang4}}\\
& {\tiny \url{https://www.youtube.com/watch?v=SUTVS1B8Kag&t=139s}}
\\
\end{tabular}
\end{table}

\end{frame}

\begin{frame}[fragile]{Suprafixes}
\begin{table}
\begin{tabular}{l||c >{\centering\arraybackslash}p{0.45cm} >{\centering\arraybackslash}p{0.45cm} c||c c}
    \hline
     & *A & \multicolumn{2}{|c|}{*B} & *C & \multicolumn{2}{c}{*D}\\
     & & & & & {\scriptsize L} & \multicolumn{1}{|c}{{\scriptsize S}}\\ \hline
    陰 & 1 ˥(\~{}˥˧) & \multicolumn{2}{|c|}{2 ˨˥\~{}˧˥} & 3 ˧ & 3 ˧ & \multicolumn{1}{|c}{1 ˥}\\ \hline 
    陽 & 4 ˨˩\~{}˩ & \multicolumn{2}{|c|}{5 ˩˧} & 6 ˨ & \multicolumn{2}{c}{6 ˨}\\ \hline
\end{tabular}
\end{table}
\medskip
Quiz for Cantonese speakers:\\
What is the citation tone of these syllables pronounced \textbf{jyun2}?\\
\begin{table}
    \flushleft
    \begin{tabular}{r l l l l}
    & & Cantonese & \textcolor{gray!75}{Mandarin} & \textcolor{gray!75}{Sino-Viet}\\
    garden & 園 & jyun? & \textcolor{gray!75}{yuán} & \textcolor{gray!75}{viên / vườn}\\
    pellet & 丸 & jyun? & \textcolor{gray!75}{wán} & \textcolor{gray!75}{hoàn}\\
    courtyard & 院 & jyun? & \textcolor{gray!75}{yuàn} & \textcolor{gray!75}{viện}\\
    (surname) & 阮 & jyun? & \textcolor{gray!75}{ruǎn} & \textcolor{gray!75}{nguyễn}\\
    garden & 苑 & jyun? & \textcolor{gray!75}{yuàn} & \textcolor{gray!75}{uyển / uốn}\\ 
    \end{tabular}
\end{table}

\end{frame}

\subsection{Suprafixes 2}

\begin{frame}[fragile]{Suprafixes}
\textbf{On Tone 1 [˥] [˥˧] again:}\\
\medskip
Note: Modern speakers still have two words which must be [˥˧] (or maybe this is boundary intonation):\\
\medskip
1.\\
先 sin1 [˥˧], sentence final particle expressing reservations in a question, vs.\\
先 sin1 [˥]\textsuperscript{usually}, adverb \enquote{first, before, earlier}\\
\medskip
\begin{table}
    \raggedright
    \begin{tabular}{l l l l l l l l}
        係& 唔& 係& 你& 去& 先& 先& ？\\
        hai6& m4& hai6& nei5& heoi3& sin˥& sin˥˧& ?\\
        be& \textsc{neg}& be& \textsc{2sg}& go& first& \textsc{sfp}& \\
        \multicolumn{8}{l}{\enquote*{Is it that you are going first...?}}
    \end{tabular}
\end{table}

\end{frame}

\begin{frame}[fragile]{Suprafixes}
2.\\
唓 ce1 [˥˧] interjection expressing dismissal, vs.\\
車 ce1 [˥]\textsuperscript{usually} \enquote*{vehicle}\\
\bigskip
Otherwise, nowadays [˥] [˥˧] are free variations for most speakers.\\
e.g. 根 gan1 [˥]\~{}[˥˧] \enquote*{root}\\
Hong Kong: tone 1 most usually [˥]\\
\medskip
%\emph{\footnotesize (Note: free varation for sonorant-ending syllables;\\
%obstruent-ending syllable tone 1 always ˥\\
%e.g. 拮 gat1 [˥] \enquote*{to prick})}\\
\end{frame}

\begin{frame}[fragile]{Suprafixes}
For people who still make a distinction between /˥˧/ and /˥/ {\scriptsize (Cheung 1969)}:\\ 
\smallskip
Tone 1 [˥˧] vs. Tone *1 [˥]\\ 
Suprafix tone *1: primarily on nouns, can be applied to morphemes in any tone \textsuperscript{other than tone 7 [˥]}, e.g.:\\
\begin{table}
\flushleft
\begin{tabular}{l l}
    忠 zung1 \enquote*{loyal(ty)} & 鐘 zung1*1 \enquote*{clock}\\
    好酸 hou2 syun1 & 好孫 hou2 syun1*1\\
    \quad \enquote*{very sour} & \quad \enquote*{good grandchild}\\
    四方 sei3 fong1 & 四方 sei3 fong1*1\\
    \quad \enquote*{four directions} & \quad \enquote*{square}\\
    閂門 saan1 mun4 & 門閂 mun4 saan1*1\\
    \quad \enquote*{close door} & \quad \enquote*{latch}
\end{tabular}
\end{table}
Female given names can end in tone *1 (\textsc{dim}) but not tone 1\\ Male name varies
\vfill \setstretch{0.65} {\tiny Cheung, Yat-Shing 張日昇. 1969. 香港粵語陰平調及變調問題.  香港中文大學中國文化研究所學報 2: 81–107. \url{https://www.cuhk.edu.hk/ics/journal/articles/v02p081.pdf}}\par 
\end{frame}

\begin{frame}[fragile]{Suprafixes}
Tone 2:\\
Cheung 張 (1969: 96-97): only one tone 2 in Hong Kong\\
Gāo 高 (1980: 7): Tone *2 rises higher than Tone 2 in Canton\\
Chao (1947: 34-35): Tone 2 [˧˥], Tone *2 [˨˥]\\
Ball (1888: xxxi) \enquote{third rising tone} [Tone *2]: \enquote{banning lower and ascending higher than [the other two rising tones]}, something like [˧˩˥]\\
\begin{table}
\flushleft
\begin{tabular}{l l}
    訪 fong2 \enquote*{to visit} & 房 fong4*2 \enquote*{room}\\
    抰 joeng2 \enquote*{to shake} & 樣 joeng6*2 \enquote*{pattern}\\
    嚿屎 gau6 si2 \enquote*{lump faeces} & 舊時 gau6 si4*2 \enquote*{the past}\\
\end{tabular}
\end{table}

\vfill \setstretch{0.65} {\tiny Gāo, Huánián 高華年. 1980. 廣州方言研究. Hong Kong: The Commercial
Press 商務印書館.\\
Chao, Yuen Ren. 1947. \emph{Cantonese Primer}. Cambridge and Massachusetts: Harvard University Press.\\
Ball, James Dyer. 1888. \emph{Cantonese Made Easy: A Book of Simple Sentences in the Canton Dialect, with Free and Literal Translations, and
Directions for the Rendering of English Grammatical Forms in Chinese}, 2nd ed. Hong Kong: China Mail Office.}

\end{frame}

\begin{frame}[fragile]{Suprafixes}
Canton \& Hong Kong at the forefront of the development of Cantonese suprafixes.\\
\bigskip
Other Cantonese varieties, e.g.\\
\smallskip
In 北海 Bak1 Hoi2 / Běihǎi Canto., out of 2179 lexical items:\\
\begin{quote}例子很少，在我们的词汇表里，我们一共记录了七个有变调的词。{\scriptsize (Chen \& Chen 2005: 26)} \end{quote}
\smallskip
In 南寧 Naam4 Ning4 / Nánníng Canto., the tone 2 suprafix:\\
\begin{quote}使用远无广州、香港白话那么普遍，基本可以穷尽。{\scriptsize (Lin \& Qin 2008: 90)}\end{quote}

\vfill \setstretch{0.65} {\tiny Chén, Xiǎojǐn 陈晓锦, and Tāo Chén 陈滔. 2005. 广西北海市粤方言调查研究. Beijing: China Social Science Press 中国社会科学出版社.\\
Lín, Yì 林亦, and Fèngyú Qín 覃凤余. 2008. 广西南宁白话研究.
Guilin: Guangxi Normal University Press 广西师范大学出版社.}

\end{frame}

\begin{frame}[fragile]{Suprafixes}
e.g.\\
\begin{table}
\flushleft
\begin{tabular}{l l l l}
     & 南寧 Canto. & Older 澳門 Canto. & 香港 Canto.\\
    白話 & baak6 waa6 & baak6 waa6*2 & baak6 waa6*2\\
    帽 & mu6 & mou6 & mou6*2\\
    枱 & toi4 & toi4(*2) & toi4*2\\
    枱布 & toi4 bu3 & toi4 bou3 & toi4*2 bou3\\
    乞兒 & hat1 ji4*1 & haat1 ji4*1 & hat1 ji4*1\\
    乞兒仔 & ? & haat1 ji4 zai2 & hat1 ji4*1 zai2\\
    澳門 & u3 mun4 & ou3 mun4(*2) & ou3 mun4*2\\
\end{tabular}
\end{table}
\medskip
{\small In Macau, /ou3 mun4/ is on average\\ \quad more formal than /ou3 mun4*2/\\
e.g. 澳門大學 \enquote*{Universidade de Macau} more likely pronounced\\ /ou3 mun4 daai6 hok6/ in formal speech.\\ (/ou3 mun4*2 daai6 hok6/ is also OK)}

\end{frame}

\begin{frame}[fragile]{Suprafixes}
What is the origin of the lexical suprafixes in Cantonese?\\
\medskip
{\scriptsize We don't actually know for sure.}\\
\medskip
Ball (1888) describes the suprafix rising tone as a long tone.\\
\bigskip
\emph{Presumably there was a suffix originally.}\\
\medskip
\emph{cf. grammatical suprafixes\\
sik6*2 (eat:\textsc{pfv}) < 食咗 sik6 zo2 (or 食嘵 sik6 hiu1)}\\
\medskip
Lexical suprafix: ?\\
\medskip
Best theory so far: < \textsc{dim} suffix -兒 -ɲi4*1 
\end{frame}

\begin{frame}[fragile]{Suprafixes}
兒 MC ɲe\textsuperscript{A} \enquote*{male child} later \enquote*{child}\\
\bigskip
南寧位子淥平話 (Pinghua, south central Guangxi)\\
雞兒 kɐj˥˧ ɲi˨˩ (fowl-\textsc{dim})
\quad 刀兒 taw˥˧ ɲi˨˩ (knife-\textsc{dim})\\
車兒 tʃʰɛ˥˧ ɲi˨˩ \~{} ˥˧ (vehicle-\textsc{dim})\\
鐺兒 tʃʰɛŋ˥˧ ɲi˥˧ (wok-\textsc{dim})\\
\medskip
桂平白話 (Cantonese Yue, central eastern Guangxi)\\
only -ŋi˥, e.g\\
雞兒 kɐj˥ ŋi˥\\
\medskip
信宜 (Yue, southwestern Guangdong)\\
豬 tʃy˥˧ > 豬\textsuperscript{兒} tʃyn↗ (pig:\textsc{dim})\\
豬兒 tʃy˥˧ ɲi˥˧ \enquote*{piglet}> 豬兒\textsuperscript{兒} tʃy˥˧ ɲin↗ (pig-\textsc{dim}:\textsc{dim})\\
\medskip
Further east: -兒 suffix no longer productive

\end{frame}

\begin{frame}[fragile]{Suprafixes}
Standard Cantonese, fossilised cases of -兒 ji4*1\\
\medskip
乞兒 hat1 ji4*1 \enquote*{beggar}\\
捉兒人 zuk1 ji4*1 jan4*1 \enquote*{hide and seek}\\
\textcolor{gray!75}{貓兒 maau1 ji4*1 \enquote*{kitten} (moribund)}\\
\medskip
\begin{table}
    \flushleft
    \textcolor{gray!75}{
    \begin{tabular}{l l l l l}
        老 & 狗 & 嫩 & 貓 & 兒\\
        lou5 & gau2 & nyun6 & maau1 & ji2*1\\
        食 & 死 & 冇 & 人 & 知\\
        sik6 & sei2 & mou5 & jan4 & zi1\\ 
    \end{tabular}}
\end{table}
\bigskip
so... perhaps mui6*2 \& -mui6*1 < 妹兒 mui6 ji4*1 (?)\\
\smallskip
\emph{\footnotesize (The 兒 theory is the best theory so far; see also de Sousa (2024) for criticisms of this theory)}

\end{frame}

\begin{frame}[fragile]{Suprafixes}

仔 zai2 (\textsc{dim}): another possible origin of the lexical suprafixes. 
However:\\
\begin{itemize}
    \item 仔 zai2 was in ordinary tone 2 and not suprafix tone *2\\
    \item 仔 zai2 explains tone *2 [˨˥] reasonably, but tone *1 [˥] not well\\
    {\scriptsize e.g. mui6 + zai2 > mui6*2 is OK, but mui6*1 (?)}
    \item lexical suprafixes and -仔 -zai2 already coexisted in the 19\textsuperscript{th} century Cantonese textbooks, e.g. 檯仔 toi4*2 zai2 (table \textsc{dim})\\
\end{itemize}
\smallskip
{\footnotesize However, 仔 zai2 could have inspired further suprafixed forms, e.g. 麻仔 maa4 zai2 > maa4*2 \enquote*{measles} (Kwok 2016: 307)}

\vfill \setstretch{0.65}{\tiny Kwok, Bit‑Chee 郭必之. 2016. 粵語方言小稱變音的類型及其歷史來源
\emph{Chinese Studies} 漢學研究 34: 281–31.\\
de Sousa, Hilário. 2024. Some observations on the Cantonese lexical suprafixes. \emph{Languages} 9 (10). 311. \url{https://doi.org/10.3390/languages9100311}}

\end{frame}

\end{document}
